% GENERATED FILE. Edit canonical lessons, then run npm run generate:books.
% canonical-source-hash: fnv1a64:1495fae11d61bd31
% canonical-lessons: TA-C516-ulavi, TA-C516-tetupori, TA-C516-valaippativu, TA-C516-kanoli, TA-C516-payanar

\chapter{Online}
\label{ch:ta-online}
\hlchaptermodality{516}

\begin{chapteropening}
\textbf{By the end of this chapter:} \emph{I can talk about browsing, blogs and videos.}

Complete the last lesson of chapter 516: I can talk about browsing, blogs and videos.
\end{chapteropening}

\section[\texorpdfstring{\ta{உலாவி} (ulāvi)}{உலாவி (ulāvi)}]{\ta{உலாவி} (ulāvi) — a web browser}

\label{lesson:TA-C516-ulavi}



\begin{quote}
Before the new one: say the Tamil for a laptop, then the Tamil for an e-book.
\end{quote}

\subsection*{You'll want to know: \ta{உலாவி}}
\textbf{\ta{உலாவி}} — \emph{ulāvi} — \textquotedblleft{}a web browser\textquotedblright{}.

\begin{quote}
\ta{என்} \ta{உலாவி} \ta{மிகவும்} \ta{மெதுவாக} \ta{வேலை} \ta{செய்கிறது}.
\end{quote}

\begin{quote}
\emph{eṉ ulāvi mikavum metuvāka vēlai ceykiṟatu.}
\end{quote}

\begin{quote}
My browser is working very slowly.
\end{quote}

From \ta{உலா}, a stroll: the program that strolls through websites.

\begin{cousinweb}
\textbf{\ta{உலாவி}} (\emph{ulāvi}) is a modern coinage built on the native \ta{உலா} (\emph{ulā}), \textquotedblleft{}a stroll, a round\textquotedblright{}, with -\ta{இ}, \textquotedblleft{}the one that does it\textquotedblright{}. DEDR 2693 gives \ta{உலா}, \textquotedblleft{}a procession\textquotedblright{}, beside \ta{உலாவு} (\emph{ulāvu}), \textquotedblleft{}to stroll\textquotedblright{}.
\end{cousinweb}

\begin{grammarlens}[title={What you've built}]
One more word for describing what you lived and felt.
\end{grammarlens}

\subsection*{Your turn}
\begin{itemize}
\raggedright
  \item \emph{Say it:} \emph{ulāvi}
  \item \emph{Say it:} \emph{ulāvi}, once more
  \item \emph{Say it:} \emph{eṉ ulāvi mikavum metuvāka vēlai ceykiṟatu.} — My browser is working very slowly.
  \item \emph{Recall:} say the Tamil for a laptop, then the Tamil for an e-book, then say \emph{ulāvi} again
\end{itemize}

\begin{tcolorbox}[breakable,colback=teal!4,colframe=teal!35!black,title={Before you move on}]
What does \ta{உலாவி} mean? (\textquotedblleft{}A web browser\textquotedblright{}.) Say it once more.
\end{tcolorbox}
\section[\texorpdfstring{\ta{தேடுபொறி} (tēṭupoṟi)}{தேடுபொறி (tēṭupoṟi)}]{\ta{தேடுபொறி} (tēṭupoṟi) — a search engine}

\label{lesson:TA-C516-tetupori}



\begin{quote}
Before the new one: say the Tamil for an e-book, then the Tamil for a web browser.
\end{quote}

\subsection*{You'll want to know: \ta{தேடுபொறி}}
\textbf{\ta{தேடுபொறி}} — \emph{tēṭupoṟi} — \textquotedblleft{}a search engine\textquotedblright{}.

\begin{quote}
\ta{நான்} \ta{தேடுபொறியில்} \ta{அந்தப்} \ta{பாடலைத்} \ta{தேடினேன்}.
\end{quote}

\begin{quote}
\emph{nāṉ tēṭupoṟiyil antap pāṭalait tēṭiṉēṉ.}
\end{quote}

\begin{quote}
I looked for that song on a search engine.
\end{quote}

\ta{தேடு} \textquotedblleft{}to search\textquotedblright{} + \ta{பொறி} \textquotedblleft{}a device\textquotedblright{}.

\begin{cousinweb}
\textbf{\ta{தேடுபொறி}} (\emph{tēṭupoṟi}) is a modern compound: the native \ta{தேடு} (\emph{tēṭu}), \textquotedblleft{}to search\textquotedblright{} (DEDR 3456), and \ta{பொறி} (\emph{poṟi}), \textquotedblleft{}a machine, a device\textquotedblright{}; the origin of the device sense of \ta{பொறி} is not settled.
\end{cousinweb}

\begin{grammarlens}[title={What you've built}]
One more word for describing what you lived and felt.
\end{grammarlens}

\subsection*{Your turn}
\begin{itemize}
\raggedright
  \item \emph{Say it:} \emph{tēṭupoṟi}
  \item \emph{Say it:} \emph{tēṭupoṟi}, once more
  \item \emph{Say it:} \emph{nāṉ tēṭupoṟiyil antap pāṭalait tēṭiṉēṉ.} — I looked for that song on a search engine.
  \item \emph{Recall:} say the Tamil for an e-book, then the Tamil for a web browser, then say \emph{tēṭupoṟi} again
\end{itemize}

\begin{tcolorbox}[breakable,colback=teal!4,colframe=teal!35!black,title={Before you move on}]
What does \ta{தேடுபொறி} mean? (\textquotedblleft{}A search engine\textquotedblright{}.) Say it once more.
\end{tcolorbox}
\section[\texorpdfstring{\ta{வலைப்பதிவு} (valaippativu)}{வலைப்பதிவு (valaippativu)}]{\ta{வலைப்பதிவு} (valaippativu) — a blog}

\label{lesson:TA-C516-valaippativu}



\begin{quote}
Before the new one: say the Tamil for a web browser, then the Tamil for a search engine.
\end{quote}

\subsection*{You'll want to know: \ta{வலைப்பதிவு}}
\textbf{\ta{வலைப்பதிவு}} — \emph{valaippativu} — \textquotedblleft{}a blog\textquotedblright{}.

\begin{quote}
\ta{அவள்} \ta{சமையல்} \ta{பற்றி} \ta{ஒரு} \ta{வலைப்பதிவு} \ta{எழுதுகிறாள்}.
\end{quote}

\begin{quote}
\emph{avaḷ camaiyal paṟṟi oru valaippativu eḻutukiṟāḷ.}
\end{quote}

\begin{quote}
She writes a blog about cooking.
\end{quote}

\ta{வலை} \textquotedblleft{}a net, the web\textquotedblright{} + \ta{பதிவு} \textquotedblleft{}a record\textquotedblright{}. The writer of one is a \ta{வலைப்பதிவர்}.

\begin{cousinweb}
\textbf{\ta{வலைப்பதிவு}} (\emph{valaippativu}) is built from native parts to match English \emph{weblog}: \ta{வலை} (\emph{valai}), \textquotedblleft{}a net\textquotedblright{} (DEDR 5288), now \textquotedblleft{}the web\textquotedblright{}, and \ta{பதிவு} (\emph{pativu}), \textquotedblleft{}a record, an entry\textquotedblright{} (DEDR 3911).
\end{cousinweb}

\begin{grammarlens}[title={What you've built}]
One more word for describing what you lived and felt.
\end{grammarlens}

\subsection*{Your turn}
\begin{itemize}
\raggedright
  \item \emph{Say it:} \emph{valaippativu}
  \item \emph{Say it:} \emph{valaippativu}, once more
  \item \emph{Say it:} \emph{avaḷ camaiyal paṟṟi oru valaippativu eḻutukiṟāḷ.} — She writes a blog about cooking.
  \item \emph{Recall:} say the Tamil for a web browser, then the Tamil for a search engine, then say \emph{valaippativu} again
\end{itemize}

\begin{tcolorbox}[breakable,colback=teal!4,colframe=teal!35!black,title={Before you move on}]
What does \ta{வலைப்பதிவு} mean? (\textquotedblleft{}A blog\textquotedblright{}.) Say it once more.
\end{tcolorbox}
\section[\texorpdfstring{\ta{காணொளி} (kāṇoḷi)}{காணொளி (kāṇoḷi)}]{\ta{காணொளி} (kāṇoḷi) — a video}

\label{lesson:TA-C516-kanoli}



\begin{quote}
Before the new one: say the Tamil for a search engine, then the Tamil for a blog.
\end{quote}

\subsection*{You'll want to know: \ta{காணொளி}}
\textbf{\ta{காணொளி}} — \emph{kāṇoḷi} — \textquotedblleft{}a video\textquotedblright{}.

\begin{quote}
\ta{அவன்} \ta{நாய்க்குட்டியின்} \ta{காணொளியை} \ta{எனக்கு} \ta{அனுப்பினான்}.
\end{quote}

\begin{quote}
\emph{avaṉ nāykkuṭṭiyiṉ kāṇoḷiyai eṉakku aṉuppiṉāṉ.}
\end{quote}

\begin{quote}
He sent me a video of the puppy.
\end{quote}

Plural: \ta{காணொளிகள்}, \emph{kāṇoḷikaḷ}.

\ta{காண்} \textquotedblleft{}to see\textquotedblright{} + \ta{ஒளி} \textquotedblleft{}light\textquotedblright{}: seen light, a moving picture.

\begin{cousinweb}
\textbf{\ta{காணொளி}} (\emph{kāṇoḷi}) is a modern native compound: \ta{காண்} (\emph{kāṇ}), \textquotedblleft{}to see\textquotedblright{} (DEDR 1443), and \ta{ஒளி} (\emph{oḷi}), \textquotedblleft{}light\textquotedblright{} (DEDR 1016).
\end{cousinweb}

\begin{grammarlens}[title={What you've built}]
One more word for describing what you lived and felt.
\end{grammarlens}

\subsection*{Your turn}
\begin{itemize}
\raggedright
  \item \emph{Say it:} \emph{kāṇoḷi}
  \item \emph{Say it:} \emph{kāṇoḷi}, once more
  \item \emph{Say it:} \emph{avaṉ nāykkuṭṭiyiṉ kāṇoḷiyai eṉakku aṉuppiṉāṉ.} — He sent me a video of the puppy.
  \item \emph{Recall:} say the Tamil for a search engine, then the Tamil for a blog, then say \emph{kāṇoḷi} again
\end{itemize}

\begin{tcolorbox}[breakable,colback=teal!4,colframe=teal!35!black,title={Before you move on}]
What does \ta{காணொளி} mean? (\textquotedblleft{}A video\textquotedblright{}.) Say it once more.
\end{tcolorbox}
\section[\texorpdfstring{\ta{பயனர்} (payaṉar)}{பயனர் (payaṉar)}]{\ta{பயனர்} (payaṉar) — a user}

\label{lesson:TA-C516-payanar}



\begin{quote}
Before the new one: say the Tamil for a blog, then the Tamil for a video.
\end{quote}

\subsection*{You'll want to know: \ta{பயனர்}}
\textbf{\ta{பயனர்}} — \emph{payaṉar} — \textquotedblleft{}a user\textquotedblright{}.

\begin{quote}
\ta{இந்தச்} \ta{செயலிக்கு} \ta{லட்சக்கணக்கான} \ta{பயனர்கள்} \ta{இருக்கிறார்கள்}.
\end{quote}

\begin{quote}
\emph{intac ceyalikku laṭcakkaṇakkāṉa payaṉarkaḷ irukkiṟārkaḷ.}
\end{quote}

\begin{quote}
This app has lakhs of users.
\end{quote}

Plural: \ta{பயனர்கள்}, \emph{payaṉarkaḷ}.

\ta{பயன்} \textquotedblleft{}use, benefit\textquotedblright{} + -\ta{அர்} \textquotedblleft{}person\textquotedblright{}. The plural is \ta{பயனர்கள்}.

\begin{cousinweb}
\textbf{\ta{பயனர்}} (\emph{payaṉar}) is a modern coinage: \ta{பயன்} (\emph{payaṉ}), \textquotedblleft{}use, benefit\textquotedblright{}, with the personal ending -\ta{அர்}; the origin of \ta{பயன்} is not settled.
\end{cousinweb}

\begin{grammarlens}[title={What you've built}]
One more word for describing what you lived and felt.
\end{grammarlens}

\subsection*{Your turn}
\begin{itemize}
\raggedright
  \item \emph{Say it:} \emph{payaṉar}
  \item \emph{Say it:} \emph{payaṉar}, once more
  \item \emph{Say it:} \emph{intac ceyalikku laṭcakkaṇakkāṉa payaṉarkaḷ irukkiṟārkaḷ.} — This app has lakhs of users.
  \item \emph{Recall:} say the Tamil for a blog, then the Tamil for a video, then say \emph{payaṉar} again
\end{itemize}

\begin{tcolorbox}[breakable,colback=teal!4,colframe=teal!35!black,title={Before you move on}]
What does \ta{பயனர்} mean? (\textquotedblleft{}A user\textquotedblright{}.) Say it once more.
\end{tcolorbox}
